% !TEX root = ../../main_power.tex
% ---------------------------------------------------------------------
%  Hardware-Massnahmen (Teil 3 externe Verbraucher, Teil 4 Hardware)
% ---------------------------------------------------------------------
\section{Hardware}
\label{sec:pwrhw}

\subsection{Power Gating und Clock Gating}
\label{rez:gating}
\begin{ztab}{@{}lX@{}}
\T{Power Gating} & Speisung \emph{unterbrechen} -- Power-Management-IC,
  FET, High-Side-/Low-Side-Schalter. Spart alles, aber:
  Einstellungen gehen verloren, Wiederanlauf dauert \\
\T{Clock Gating} & nur den \emph{Takt} anhalten. Spart den dynamischen
  Anteil; Spannung bleibt, \T{Einstellungen bleiben erhalten}, schnell
  wieder einschaltbar \\
\end{ztab}
Externe ICs haben dafür oft einen \R{Enable}- oder \R{Sleep}-Pin.
\begin{falle}[Parasitäre Speisung]
Ein abgeschalteter IC kann sich über seine \emph{Signalleitungen}
speisen: UART-, SPI- oder I\textsuperscript{2}C-Pegel treiben über die
internen Schutzdioden Strom in die abgeschaltete Versorgung. Der Baustein
zieht dann Strom, obwohl er „aus“ ist -- und läuft womöglich halb an.
Beim Abschalten müssen die Signalleitungen mit abgeschaltet oder auf
Eingang gelegt werden.
\end{falle}

\subsection{Bauteilwahl und Umgebung}
\begin{itemize}
\item Komponenten wählen, die \emph{auf Low Power optimiert} sind
\item \T{3.3\,V statt 5.0\,V} Logikspannung -- aber Achtung Störfestigkeit
\item ICs mit Shutdown- oder Sleep-Modus
\item hochwertige Speisung bzw.\ Spannungswandler (Effizienz, Up/Down)
\item ICs, die keine Kühlung brauchen; Schutz vor Umgebungswärme
\item Kondensatoren: Leckströme (Elektrolyt), Alterung
\item unbenutzte Pins: Datenblatt lesen -- definierter Eingang, nicht
      treiben, internes Pull-up je nach Hersteller
\end{itemize}
\SW{Anzeigen} der oft unterschätzte Verbraucher: Low-Power-LCD statt
LED/OLED, Fläche und Helligkeit reduzieren, Bildschirmaktivität
vermindern, dimmen bzw.\ modulieren, im Idealfall \T{bi-stabile
(Zero-Power-) Anzeigen}, die nur beim Umschalten Strom brauchen.
\begin{bsp}[Low Power kostet Geld]
Zwei funktionsgleiche ICs aus dem Katalog: TSSOP-16 mit 1\,mA für
72.80 (100 Stück) gegen TSSOP-20 mit 1\,\textmu A für 104.00. Faktor
1000 beim Strom, Faktor 1.4 beim Preis -- die Abwägung gehört in die
Argumentation.
\end{bsp}

\subsection{Halbleiterfläche}
\label{sec:flaeche}
\begin{formel}[Fläche ist der gemeinsame Nenner]
Fläche $\propto$ Kosten (weniger ICs pro Wafer) \emph{und}
Fläche $\propto C$ \ra dynamischer Verbrauch. Jede eingesparte
Funktionseinheit spart beides.
\end{formel}
\begin{ztab}{@{}lX@{}}
Speicher & RAM und Cache sind gross und brauchen viel Energie \ra
  Alternativen MRAM, FRAM. Flash ist dichter und sparsamer, braucht aber
  $V_{prog}$. Nur so viel Speicher wie nötig -- schwierig wegen Stack,
  späteren Funktionen, Produktlebenszyklus und Entwicklungssupport
  (Logs, Trace) \\
Takt & \R{ICG} (Internal Clock Generator) gegen externen Takt;
  \R{PMC} (Power Management Controller) mit Ladungspumpen für
  $V_{FlashProg}$ \\
Core & Instruktionen, Register, Pipelines, Caches \\
Peripherie & nur so viel wie nötig; der \R{ADC} braucht auffällig
  viel Fläche \\
\end{ztab}
